\documentclass[11pt,a4paper]{article}

% =========================
% Packages
% =========================
\usepackage[utf8]{inputenc}
\usepackage[T1]{fontenc}
\usepackage[french]{babel}
\usepackage[a4paper,margin=2.2cm]{geometry}
\usepackage{lmodern}
\usepackage{amsmath,amssymb,amsthm,mathtools}
\usepackage{enumitem}
\usepackage{xcolor}
\usepackage{fancyhdr}
\usepackage{lastpage}
\usepackage{hyperref}
\usepackage{array}
\usepackage{stmaryrd}

% =========================
% Hyperref
% =========================
\hypersetup{
    colorlinks=true,
    linkcolor=black,
    urlcolor=blue,
    pdftitle={Corrigé détaillé - DM de géométrie},
    pdfauthor={},
    pdfsubject={Corrigé détaillé de géométrie analytique}
}

% =========================
% General layout
% =========================
\setlength{\parindent}{0pt}
\setlength{\parskip}{0.5em}

\pagestyle{fancy}
\fancyhf{}
\fancyhead[L]{\textsc{Corrigé détaillé}}
\fancyhead[C]{\textsc{Géométrie analytique}}
\fancyhead[R]{\textsc{Plan et espace}}
\fancyfoot[C]{\thepage/\pageref{LastPage}}

% =========================
% Custom commands
% =========================
\newcommand{\R}{\mathbb{R}}
\newcommand{\vv}[1]{\overrightarrow{#1}}
\newcommand{\Boule}{\mathcal{B}}
\newcommand{\Sphere}{\mathcal{S}}
\newcommand{\Cube}{\mathcal{C}}
\newcommand{\dd}{\,\mathrm d}
\newcommand{\coord}[1]{\left(#1\right)}
\newcommand{\norme}[1]{\left\lVert #1\right\rVert}

\newtheorem*{methode}{Méthode}
\newtheorem*{idee}{Idée}
\newtheorem*{remarque}{Remarque}

% =========================
% Document
% =========================
\begin{document}

\begin{center}
    {\Large \textbf{Corrigé détaillé --- Devoir Maison de géométrie analytique}}\\[0.4em]
    {\large \textbf{Raisonnement dans le plan et dans l'espace}}\\[0.8em]
    \rule{0.75\textwidth}{0.6pt}
\end{center}

\vspace{0.8em}

\section*{Problème 1 --- Boule dans l'espace, intersections et inclusion}

On considère
\[
\Omega(1,1,1),\qquad 
\Boule=\left\{M(x,y,z)\in\R^3\mid \Omega M\le 3\right\},
\]
et sa frontière
\[
\Sphere=\left\{M(x,y,z)\in\R^3\mid \Omega M=3\right\}.
\]
On considère aussi le cube
\[
\Cube=[0,2]\times[0,2]\times[0,2].
\]

\subsection*{Partie A --- Équation, appartenance, inclusion}

\begin{enumerate}[label=\textbf{\arabic*.}]
    \item \textbf{Équation cartésienne de \(\Sphere\).}

    Par définition,
    \[
    M(x,y,z)\in \Sphere \iff \Omega M=3.
    \]
    Comme le repère est orthonormé,
    \[
    \Omega M^2=(x-1)^2+(y-1)^2+(z-1)^2.
    \]
    Donc
    \[
    M\in \Sphere \iff (x-1)^2+(y-1)^2+(z-1)^2=9.
    \]

    Ainsi, une équation cartésienne de \(\Sphere\) est
    \[
    \boxed{(x-1)^2+(y-1)^2+(z-1)^2=9.}
    \]

    \item \textbf{Appartenance de \(A\), \(B\), \(C\) à \(\Sphere\).}

    On calcule :

    \[
    A(4,1,1)\quad\Rightarrow\quad (4-1)^2+(1-1)^2+(1-1)^2=3^2=9.
    \]
    Donc \(A\in \Sphere\).

    \[
    B(1,4,1)\quad\Rightarrow\quad (1-1)^2+(4-1)^2+(1-1)^2=3^2=9.
    \]
    Donc \(B\in \Sphere\).

    \[
    C(1,1,-2)\quad\Rightarrow\quad (1-1)^2+(1-1)^2+(-2-1)^2= (-3)^2=9.
    \]
    Donc \(C\in \Sphere\).

    Finalement,
    \[
    \boxed{A,B,C\in \Sphere.}
    \]

    \item \textbf{Le point \(D(2,2,2)\) appartient-il à \(\Boule\) ?}

    On calcule :
    \[
    \Omega D^2=(2-1)^2+(2-1)^2+(2-1)^2=1+1+1=3.
    \]
    Ainsi
    \[
    \Omega D=\sqrt3<3.
    \]
    Donc \(D\) est à distance strictement inférieure à \(3\) du centre \(\Omega\), d'où
    \[
    \boxed{D\in \Boule.}
    \]

    \item \textbf{Montrer que \(\Cube\subset \Boule\).}

    Soit \(M(x,y,z)\in \Cube\). Alors
    \[
    0\le x\le 2,\qquad 0\le y\le 2,\qquad 0\le z\le 2.
    \]
    Donc
    \[
    -1\le x-1\le 1,\qquad -1\le y-1\le 1,\qquad -1\le z-1\le 1.
    \]
    En particulier,
    \[
    (x-1)^2\le 1,\qquad (y-1)^2\le 1,\qquad (z-1)^2\le 1.
    \]
    En sommant,
    \[
    \Omega M^2=(x-1)^2+(y-1)^2+(z-1)^2\le 3.
    \]
    Donc
    \[
    \Omega M\le \sqrt3<3.
    \]
    Ainsi \(M\in \Boule\).

    Comme ceci vaut pour tout \(M\in \Cube\), on obtient
    \[
    \boxed{\Cube\subset \Boule.}
    \]

    \item \textbf{Cette inclusion est-elle stricte ?}

    Oui. Il suffit d'exhiber un point de \(\Boule\) qui n'appartient pas à \(\Cube\).  
    Par exemple,
    \[
    E(4,1,1).
    \]
    On a vu que \(E=A\in \Sphere\), donc \(E\in \Boule\). En revanche, \(x_E=4\notin [0,2]\), donc \(E\notin \Cube\).

    Ainsi
    \[
    \boxed{\Cube\subsetneq \Boule.}
    \]
\end{enumerate}

\subsection*{Partie B --- Intersection sphère--droite}

On considère la droite \(\Delta\) passant par \(\Omega\) et de vecteur directeur
\[
\vec u=(1,2,2).
\]

\begin{enumerate}[label=\textbf{\arabic*.},resume]
    \item \textbf{Représentation paramétrique de \(\Delta\).}

    Une droite passant par \(\Omega(1,1,1)\) et de vecteur directeur \((1,2,2)\) admet pour représentation paramétrique
    \[
    \boxed{
    \Delta:\ \begin{cases}
    x=1+t\\
    y=1+2t\\
    z=1+2t
    \end{cases}\quad (t\in\R).}
    \]

    \item \textbf{Intersection \(\Delta\cap \Sphere\).}

    Soit \(M\in \Delta\). Alors
    \[
    M(1+t,1+2t,1+2t).
    \]
    Pour que \(M\in \Sphere\), il faut
    \[
    (x-1)^2+(y-1)^2+(z-1)^2=9.
    \]
    Donc
    \[
    t^2+(2t)^2+(2t)^2=9
    \]
    soit
    \[
    9t^2=9.
    \]
    Ainsi
    \[
    t^2=1 \iff t=\pm 1.
    \]
    On obtient les deux points :
    \[
    t=1 \Rightarrow M_1(2,3,3),
    \]
    \[
    t=-1 \Rightarrow M_2(0,-1,-1).
    \]

    Donc
    \[
    \boxed{\Delta\cap \Sphere=\left\{(2,3,3),(0,-1,-1)\right\}.}
    \]

    \item \textbf{Interprétation géométrique.}

    La droite \(\Delta\) passe par le centre \(\Omega\) de la sphère. Une droite passant par le centre coupe la sphère en deux points diamétralement opposés.

    Ici,
    \[
    \boxed{\Delta \text{ est sécante à la sphère et l'intersection est un diamètre.}}
    \]

    \item \textbf{Intersection de \(\Delta'\) avec \(\Sphere\).}

    La droite \(\Delta'\) est donnée par
    \[
    \begin{cases}
    x=4+t\\
    y=4\\
    z=1
    \end{cases}
    \quad (t\in\R).
    \]
    Soit \(M\in \Delta'\). Alors
    \[
    M(4+t,4,1).
    \]
    Pour que \(M\in \Sphere\), il faut
    \[
    (x-1)^2+(y-1)^2+(z-1)^2=9.
    \]
    Donc
    \[
    (4+t-1)^2+(4-1)^2+(1-1)^2=9,
    \]
    soit
    \[
    (t+3)^2+9=9.
    \]
    Ainsi
    \[
    (t+3)^2=0 \iff t=-3.
    \]
    Il y a un unique point d'intersection :
    \[
    M(1,4,1).
    \]

    Finalement
    \[
    \boxed{\Delta'\cap \Sphere=\left\{(1,4,1)\right\}.}
    \]

    \item \textbf{Cas possibles pour l'intersection d'une droite et d'une sphère.}

    Une droite et une sphère peuvent :
    \begin{itemize}
        \item ne pas se couper : intersection vide ;
        \item se couper en un seul point : la droite est tangente à la sphère ;
        \item se couper en deux points : la droite est sécante à la sphère.
    \end{itemize}

    Donc les trois cas possibles sont :
    \[
    \boxed{\varnothing,\quad \{P\},\quad \{P,Q\}.}
    \]
\end{enumerate}

\subsection*{Partie C --- Intersection sphère--plan}

On considère le plan
\[
P:\ x+y+z=3.
\]

\begin{enumerate}[label=\textbf{\arabic*.},resume]
    \item \textbf{Vérifier que \(\Omega\in P\).}

    Pour \(\Omega(1,1,1)\),
    \[
    1+1+1=3.
    \]
    Donc
    \[
    \boxed{\Omega\in P.}
    \]

    \item \textbf{Nature de \(P\cap \Sphere\).}

    Le plan \(P\) passe par le centre \(\Omega\) de la sphère. L'intersection d'un plan passant par le centre d'une sphère est un cercle de centre ce centre.

    Donc
    \[
    \boxed{P\cap \Sphere \text{ est un cercle.}}
    \]

    \item \textbf{Centre et rayon de cette intersection.}

    Comme le plan passe par le centre \(\Omega\), le centre du cercle de section est \(\Omega\) lui-même.  
    Le rayon de ce cercle est alors égal au rayon de la sphère, soit \(3\).

    Donc
    \[
    \boxed{\text{centre } \Omega(1,1,1)\quad\text{et rayon }3.}
    \]

    \item \textbf{Distance de \(\Omega\) au plan \(Q:x+y+z=7\).}

    On écrit \(Q\) sous la forme
    \[
    x+y+z-7=0.
    \]
    La distance d'un point \((x_0,y_0,z_0)\) au plan \(ax+by+cz+d=0\) vaut
    \[
    \frac{|ax_0+by_0+cz_0+d|}{\sqrt{a^2+b^2+c^2}}.
    \]
    Ici,
    \[
    d(\Omega,Q)=\frac{|1+1+1-7|}{\sqrt{1^2+1^2+1^2}}
    =\frac{4}{\sqrt3}.
    \]
    Donc
    \[
    \boxed{d(\Omega,Q)=\frac{4}{\sqrt3}.}
    \]

    \item \textbf{Nature de \(Q\cap \Sphere\).}

    Le rayon de la sphère vaut \(3\). Or
    \[
    \frac{4}{\sqrt3}<3
    \]
    car
    \[
    \frac{16}{3}<9.
    \]
    Donc le plan est à une distance strictement inférieure au rayon : il coupe la sphère suivant un cercle.

    Ainsi
    \[
    \boxed{Q\cap \Sphere \text{ est un cercle.}}
    \]

    \item \textbf{Nature de \(R\cap \Sphere\) avec \(R:x+y+z=10\).}

    On calcule
    \[
    d(\Omega,R)=\frac{|1+1+1-10|}{\sqrt3}=\frac{7}{\sqrt3}.
    \]
    Or
    \[
    \frac{7}{\sqrt3}>3
    \]
    car
    \[
    \frac{49}{3}>9.
    \]
    Le plan est donc trop loin du centre pour rencontrer la sphère.

    Donc
    \[
    \boxed{R\cap \Sphere=\varnothing.}
    \]

    \item \textbf{Résumé général.}

    Si une sphère de centre \(\Omega\) et de rayon \(r\) rencontre un plan \(\Pi\), alors :
    \begin{itemize}
        \item si \(d(\Omega,\Pi)>r\), alors \(\Pi\cap \Sphere=\varnothing\) ;
        \item si \(d(\Omega,\Pi)=r\), alors \(\Pi\cap \Sphere\) est réduit à un point : le plan est tangent ;
        \item si \(d(\Omega,\Pi)<r\), alors \(\Pi\cap \Sphere\) est un cercle.
    \end{itemize}

    Ainsi
    \[
    \boxed{
    \begin{cases}
    d>r &\Rightarrow \varnothing,\\
    d=r &\Rightarrow \text{un point},\\
    d<r &\Rightarrow \text{un cercle}.
    \end{cases}}
    \]
\end{enumerate}

\subsection*{Partie D --- Intersection plan--cube}

On considère le plan
\[
H:\ x+y+z=2.
\]

\begin{enumerate}[label=\textbf{\arabic*.},resume]
    \item \textbf{Montrer que \(H\) coupe le cube \(\Cube\).}

    Le point \((2,0,0)\) appartient au cube car ses coordonnées appartiennent à \([0,2]\). De plus,
    \[
    2+0+0=2.
    \]
    Donc \((2,0,0)\in H\cap \Cube\).

    Ainsi
    \[
    \boxed{H\cap \Cube\neq\varnothing.}
    \]

    \item \textbf{Sommets du cube appartenant à \(H\).}

    Les sommets du cube sont les points de coordonnées \(0\) ou \(2\).  
    On cherche ceux vérifiant \(x+y+z=2\). Comme chaque coordonnée vaut \(0\) ou \(2\), la somme vaut \(2\) si et seulement si une seule coordonnée vaut \(2\) et les deux autres valent \(0\).

    Les sommets cherchés sont donc
    \[
    \boxed{(2,0,0),\ (0,2,0),\ (0,0,2).}
    \]

    \item \textbf{Points d'intersection avec les arêtes du cube.}

    Les trois points trouvés ci-dessus sont déjà des sommets, donc en particulier des points de certaines arêtes.

    Regardons s'il existe d'autres points d'intersection sur les arêtes.

    Une arête du cube est obtenue en fixant deux coordonnées à \(0\) ou \(2\), et en laissant la troisième varier dans \([0,2]\).

    \begin{itemize}
        \item Si deux coordonnées sont \(0\) et la troisième varie, l'équation \(x+y+z=2\) impose que cette troisième coordonnée vaille \(2\). On obtient un sommet.
        \item Si l'une des coordonnées fixées vaut déjà \(2\) et une autre vaut \(2\), alors la somme vaut au moins \(4\), impossible.
        \item Si une coordonnée fixée vaut \(2\) et une autre \(0\), la variable devrait alors valoir \(0\), ce qui redonne encore un sommet.
    \end{itemize}

    Il n'y a donc pas d'autre point.

    Ainsi les points d'intersection avec les arêtes sont exactement
    \[
    \boxed{(2,0,0),\ (0,2,0),\ (0,0,2).}
    \]

    \item \textbf{Nature géométrique de \(H\cap \Cube\).}

    Comme le plan coupe le cube selon trois points non alignés
    \[
    A(2,0,0),\quad B(0,2,0),\quad C(0,0,2),
    \]
    la section est le triangle \(ABC\).

    Donc
    \[
    \boxed{H\cap \Cube \text{ est le triangle de sommets }(2,0,0),(0,2,0),(0,0,2).}
    \]

    \item \textbf{Aire de cette section.}

    Calculons par exemple les longueurs :
    \[
    AB=\sqrt{(-2)^2+2^2+0^2}=\sqrt8=2\sqrt2,
    \]
    \[
    AC=\sqrt{(-2)^2+0^2+2^2}=2\sqrt2,
    \]
    \[
    BC=\sqrt{0^2+(-2)^2+2^2}=2\sqrt2.
    \]
    Le triangle est donc équilatéral de côté \(2\sqrt2\).

    L'aire d'un triangle équilatéral de côté \(a\) vaut
    \[
    \frac{\sqrt3}{4}a^2.
    \]
    Avec \(a=2\sqrt2\), on a
    \[
    \mathcal A=\frac{\sqrt3}{4}\times (2\sqrt2)^2
    =\frac{\sqrt3}{4}\times 8=2\sqrt3.
    \]

    Ainsi
    \[
    \boxed{\mathcal A(H\cap \Cube)=2\sqrt3.}
    \]
\end{enumerate}

\subsection*{Partie E --- Intersection sphère--cube}

Dans toute cette partie seulement, on considère la boule fermée
\[
\Boule'=\left\{M(x,y,z)\in\R^3 \,\middle|\, \Omega M\le \sqrt3\right\}
\]
et sa frontière
\[
\Sphere'=\left\{M(x,y,z)\in\R^3 \,\middle|\, \Omega M=\sqrt3\right\},
\]
de centre \(\Omega(1,1,1)\).

\begin{enumerate}[label=\textbf{\arabic*.},resume]
    \item \textbf{Montrer que \(\Cube\cap \Sphere'\neq\varnothing\) et que \(\Cube\subset \Boule'\).}

    Soit \(M(x,y,z)\in \Cube\). Alors
    \[
    0\le x\le 2,\qquad 0\le y\le 2,\qquad 0\le z\le 2.
    \]
    Donc
    \[
    -1\le x-1\le 1,\qquad -1\le y-1\le 1,\qquad -1\le z-1\le 1.
    \]
    En élevant au carré,
    \[
    (x-1)^2\le 1,\qquad (y-1)^2\le 1,\qquad (z-1)^2\le 1.
    \]
    En sommant,
    \[
    \Omega M^2=(x-1)^2+(y-1)^2+(z-1)^2\le 3,
    \]
    donc
    \[
    \Omega M\le \sqrt3.
    \]
    Ainsi tout point du cube appartient à la boule \(\Boule'\), d'où
    \[
    \boxed{\Cube\subset \Boule'.}
    \]

    Montrons maintenant que l'intersection avec la sphère n'est pas vide.  
    Considérons par exemple le sommet \(A(0,0,0)\) du cube. Alors
    \[
    \Omega A^2=(0-1)^2+(0-1)^2+(0-1)^2=1+1+1=3,
    \]
    donc
    \[
    \Omega A=\sqrt3.
    \]
    Ainsi
    \[
    A\in \Sphere'
    \quad\text{et}\quad
    A\in \Cube.
    \]
    Donc
    \[
    \boxed{\Cube\cap \Sphere'\neq\varnothing.}
    \]

    \item \textbf{Déterminer les sommets du cube appartenant à la sphère \(\Sphere'\).}

    Les sommets du cube \(\Cube=[0,2]^3\) sont les points dont chaque coordonnée vaut \(0\) ou \(2\).  
    Soit donc un sommet \(S(x,y,z)\) avec \(x,y,z\in\{0,2\}\).

    Alors
    \[
    (x-1)^2=(y-1)^2=(z-1)^2=1,
    \]
    puisque \(0-1=-1\) et \(2-1=1\).

    Donc
    \[
    \Omega S^2=(x-1)^2+(y-1)^2+(z-1)^2=1+1+1=3,
    \]
    d'où
    \[
    \Omega S=\sqrt3.
    \]
    Ainsi tous les sommets du cube appartiennent à la sphère.

    Donc les sommets cherchés sont
    \[
    \boxed{(0,0,0),\ (2,0,0),\ (0,2,0),\ (0,0,2),\ (2,2,0),\ (2,0,2),\ (0,2,2),\ (2,2,2).}
    \]

    \item \textbf{Déterminer les arêtes du cube qui rencontrent la sphère \(\Sphere'\).}

    Toute arête du cube relie deux sommets du cube. Or on vient de montrer que tous les sommets appartiennent à la sphère \(\Sphere'\).  
    Donc chaque arête possède au moins ses deux extrémités sur la sphère, donc rencontre la sphère.

    Ainsi
    \[
    \boxed{\text{toutes les arêtes du cube rencontrent la sphère } \Sphere'.}
    \]

    \item \textbf{Pour chacune de ces arêtes, calculer explicitement le ou les points d'intersection avec la sphère \(\Sphere'\).}

    Prenons une arête type, par exemple
    \[
    [AB]\quad\text{avec}\quad A(0,0,0),\ B(2,0,0).
    \]
    Un point de cette arête s'écrit
    \[
    M(t,0,0)\qquad \text{avec } t\in[0,2].
    \]
    La condition \(M\in\Sphere'\) est
    \[
    (t-1)^2+(0-1)^2+(0-1)^2=3,
    \]
    soit
    \[
    (t-1)^2+2=3,
    \]
    donc
    \[
    (t-1)^2=1.
    \]
    Ainsi
    \[
    t=0\quad \text{ou}\quad t=2.
    \]
    Les seuls points d'intersection sont donc les extrémités de l'arête.

    Le même raisonnement vaut pour n'importe quelle autre arête : sur une arête, deux coordonnées sont fixées égales à \(0\) ou \(2\), donc leurs écarts au centre valent toujours \(\pm1\), et la troisième coordonnée varie. L'équation de la sphère impose alors que cette troisième coordonnée soit égale à \(0\) ou \(2\), c'est-à-dire que l'on soit à l'une des extrémités.

    Donc, pour chaque arête, les seuls points d'intersection avec \(\Sphere'\) sont ses deux sommets.

    \[
    \boxed{\text{Chaque arête coupe la sphère exactement en ses deux extrémités.}}
    \]

    \item \textbf{Décrire l'ensemble \(\Cube\cap\Sphere'\).}

    Tout point du cube appartient à \(\Boule'\), mais pour appartenir à la sphère \(\Sphere'\), il faut être à distance exactement \(\sqrt3\) du centre \(\Omega(1,1,1)\).

    Or, dans le cube, la distance au centre est maximale précisément aux sommets.  
    On en déduit que l'intersection du cube et de la sphère est constituée exactement des huit sommets du cube.

    Ainsi
    \[
    \boxed{\Cube\cap\Sphere'=\{\text{les huit sommets du cube}\}.}
    \]
\end{enumerate}

\vspace{1em}
\hrule
\vspace{1em}

\section*{Problème 2 --- Lieu géométrique dans le plan et structure cachée}

On considère dans un repère orthonormé du plan
\[
A(-1,2),\qquad B(5,0),\qquad C(3,4).
\]

\begin{enumerate}[label=\textbf{\arabic*.}]
    \item \textbf{Calculer \(AB\), \(AC\), \(BC\).}

    \[
    AB^2=(5-(-1))^2+(0-2)^2=6^2+(-2)^2=36+4=40,
    \]
    donc
    \[
    AB=2\sqrt{10}.
    \]

    \[
    AC^2=(3-(-1))^2+(4-2)^2=4^2+2^2=16+4=20,
    \]
    donc
    \[
    AC=2\sqrt5.
    \]

    \[
    BC^2=(3-5)^2+(4-0)^2=(-2)^2+4^2=4+16=20,
    \]
    donc
    \[
    BC=2\sqrt5.
    \]

    Finalement
    \[
    \boxed{AB=2\sqrt{10},\qquad AC=BC=2\sqrt5.}
    \]

    \item \textbf{Le triangle \(ABC\) est-il rectangle ? isocèle ?}

    On vient de trouver
    \[
    AC=BC.
    \]
    Donc le triangle \(ABC\) est isocèle en \(C\).

    Vérifions s'il est rectangle. On compare :
    \[
    AC^2+BC^2=20+20=40=AB^2.
    \]
    D'après la réciproque du théorème de Pythagore, le triangle est rectangle en \(C\).

    Ainsi
    \[
    \boxed{ABC \text{ est rectangle isocèle en } C.}
    \]

    \item \textbf{Équation de la médiatrice de \([AB]\).}

    Le milieu de \([AB]\) est
    \[
    I_{AB}\left(\frac{-1+5}{2},\frac{2+0}{2}\right)=(2,1).
    \]
    Le vecteur
    \[
    \vv{AB}=(6,-2)=2(3,-1).
    \]
    Une droite perpendiculaire à \((AB)\) a donc pour vecteur directeur \((1,3)\), car
    \[
    (3,-1)\cdot(1,3)=3-3=0.
    \]
    La médiatrice passe par \((2,1)\), donc son équation est
    \[
    y-1=3(x-2).
    \]
    On développe :
    \[
    y=3x-5.
    \]
    Donc
    \[
    \boxed{\text{Médiatrice de }[AB] :\ y=3x-5.}
    \]

    \item \textbf{Équation de la droite \((BC)\).}

    Un vecteur directeur de \((BC)\) est
    \[
    \vv{BC}=(3-5,4-0)=(-2,4)=(-1,2).
    \]
    Donc la pente est \(-2\). Comme la droite passe par \(B(5,0)\),
    \[
    y=-2(x-5).
    \]
    Ainsi
    \[
    y=-2x+10.
    \]
    Donc
    \[
    \boxed{(BC):\ y=-2x+10.}
    \]

    \item \textbf{Intersection de la médiatrice de \([AB]\) et de \((BC)\).}

    On résout
    \[
    3x-5=-2x+10.
    \]
    Donc
    \[
    5x=15 \iff x=3.
    \]
    Puis
    \[
    y=3\times 3-5=4.
    \]
    Le point d'intersection est donc
    \[
    I(3,4).
    \]

    Ainsi
    \[
    \boxed{I=(3,4).}
    \]

    On remarque que c'est précisément le point \(C\).

    \item \textbf{Interprétation géométrique.}

    Comme \(I\) appartient à la médiatrice de \([AB]\), il est équidistant de \(A\) et de \(B\). Donc
    \[
    IA=IB.
    \]
    Ici \(I=C\), donc cela signifie
    \[
    CA=CB.
    \]
    On retrouve bien le fait que le triangle est isocèle en \(C\).

    \item \textbf{Montrer que \(\vv{MA}\cdot\vv{MB}=0\) définit un cercle.}

    Soit \(M(x,y)\). Alors
    \[
    \vv{MA}=(-1-x,\,2-y),\qquad \vv{MB}=(5-x,\,-y).
    \]
    La condition
    \[
    \vv{MA}\cdot\vv{MB}=0
    \]
    devient
    \[
    (-1-x)(5-x)+(2-y)(-y)=0.
    \]
    Développons :
    \[
    (-1-x)(5-x)=x^2-4x-5,
    \]
    et
    \[
    (2-y)(-y)=y^2-2y.
    \]
    Donc
    \[
    x^2-4x-5+y^2-2y=0,
    \]
    soit
    \[
    x^2-4x+y^2-2y=5.
    \]
    On complète les carrés :
    \[
    (x-2)^2-4+(y-1)^2-1=5,
    \]
    donc
    \[
    (x-2)^2+(y-1)^2=10.
    \]

    C'est l'équation d'un cercle.

    Ainsi
    \[
    \boxed{\vv{MA}\cdot\vv{MB}=0 \iff (x-2)^2+(y-1)^2=10.}
    \]

    \item \textbf{Centre et rayon de ce cercle.}

    L'équation
    \[
    (x-2)^2+(y-1)^2=10
    \]
    montre que le cercle a pour centre
    \[
    \boxed{(2,1)}
    \]
    et pour rayon
    \[
    \boxed{\sqrt{10}.}
    \]

    \item \textbf{Interprétation géométrique sur l'angle \(\widehat{AMB}\).}

    La condition
    \[
    \vv{MA}\cdot\vv{MB}=0
    \]
    signifie que les vecteurs \(\vv{MA}\) et \(\vv{MB}\) sont perpendiculaires. Donc l'angle
    \[
    \widehat{AMB}
    \]
    est droit.

    On retrouve donc le théorème classique :
    \[
    \boxed{\widehat{AMB}=90^\circ \iff M \text{ appartient au cercle de diamètre } [AB].}
    \]

    Vérifions la cohérence : le milieu de \([AB]\) est \((2,1)\), et
    \[
    AB=2\sqrt{10},
    \]
    donc le rayon du cercle de diamètre \([AB]\) est bien
    \[
    \frac{AB}{2}=\sqrt{10}.
    \]
    C'est exactement le cercle trouvé.

    \item \textbf{Points de \((AC)\) vérifiant \(\vv{MA}\cdot\vv{MB}=0\).}

    La droite \((AC)\) passe par \(A(-1,2)\) et \(C(3,4)\).  
    Un vecteur directeur est
    \[
    \vv{AC}=(4,2)=2(2,1),
    \]
    donc une représentation paramétrique est
    \[
    M(-1+2t,\ 2+t).
    \]
    On remplace dans l'équation du cercle :
    \[
    (x-2)^2+(y-1)^2=10.
    \]
    Cela donne
    \[
    (-1+2t-2)^2+(2+t-1)^2=10,
    \]
    soit
    \[
    (2t-3)^2+(t+1)^2=10.
    \]
    Développons :
    \[
    4t^2-12t+9+t^2+2t+1=10
    \]
    donc
    \[
    5t^2-10t+10=10
    \]
    puis
    \[
    5t^2-10t=0
    \]
    soit
    \[
    5t(t-2)=0.
    \]
    Donc
    \[
    t=0 \quad \text{ou}\quad t=2.
    \]
    Pour \(t=0\), on obtient
    \[
    M=A(-1,2).
    \]
    Pour \(t=2\), on obtient
    \[
    M(3,4)=C.
    \]

    Ainsi
    \[
    \boxed{\text{les points cherchés sont }A(-1,2)\text{ et }C(3,4).}
    \]

    \begin{remarque}
    Le point \(A\) vérifie bien l'équation algébrique, car
    \[
    \vv{AA}\cdot\vv{AB}=0.
    \]
    Géométriquement, l'angle en \(A\) n'est pas un angle ordinaire du triangle \(AMB\), mais algébriquement la condition reste vraie.
    \end{remarque}

    \item \textbf{Montrer que \(\Gamma=\{M\mid MA^2-MB^2=10\}\) est une droite.}

    On écrit
    \[
    MA^2=(x+1)^2+(y-2)^2,
    \qquad
    MB^2=(x-5)^2+y^2.
    \]
    La condition
    \[
    MA^2-MB^2=10
    \]
    devient
    \[
    (x+1)^2+(y-2)^2-(x-5)^2-y^2=10.
    \]
    Développons :
    \[
    (x^2+2x+1)+(y^2-4y+4)-(x^2-10x+25)-y^2=10.
    \]
    Les termes en \(x^2\) et \(y^2\) s'éliminent, il reste
    \[
    12x-4y-20=10.
    \]
    Donc
    \[
    12x-4y=30.
    \]
    En divisant par \(2\),
    \[
    6x-2y=15.
    \]
    On peut aussi écrire
    \[
    y=3x-\frac{15}{2}.
    \]

    Ainsi
    \[
    \boxed{\Gamma \text{ est la droite d'équation } 6x-2y=15.}
    \]

    \item \textbf{Interprétation géométrique.}

    Le lieu des points tels que la différence \(MA^2-MB^2\) soit constante est une droite.  
    Ici ce n'est pas la médiatrice de \([AB]\), mais une droite parallèle à cette médiatrice.

    En effet, la médiatrice trouvée plus haut a pour équation
    \[
    y=3x-5,
    \]
    alors que \(\Gamma\) a pour équation
    \[
    y=3x-\frac{15}{2}.
    \]
    Les coefficients directeurs sont égaux.

    Donc
    \[
    \boxed{\Gamma \text{ est une droite parallèle à la médiatrice de }[AB].}
    \]

    \item \textbf{Étude de \(\Sigma=\{M\mid MA^2+MB^2=26\}\).}

    On écrit
    \[
    MA^2=(x+1)^2+(y-2)^2,\qquad MB^2=(x-5)^2+y^2.
    \]
    Donc
    \[
    MA^2+MB^2=(x+1)^2+(y-2)^2+(x-5)^2+y^2.
    \]
    Développons :
    \[
    (x^2+2x+1)+(y^2-4y+4)+(x^2-10x+25)+y^2=26.
    \]
    Soit
    \[
    2x^2-8x+2y^2-4y+30=26.
    \]
    Donc
    \[
    2x^2-8x+2y^2-4y+4=0.
    \]
    En divisant par \(2\),
    \[
    x^2-4x+y^2-2y+2=0.
    \]
    On complète les carrés :
    \[
    (x-2)^2-4+(y-1)^2-1+2=0,
    \]
    donc
    \[
    (x-2)^2+(y-1)^2=3.
    \]

    Ainsi \(\Sigma\) est un cercle de centre \((2,1)\) et de rayon \(\sqrt3\).

    Finalement
    \[
    \boxed{\Sigma=\left\{M(x,y)\in\R^2\mid (x-2)^2+(y-1)^2=3\right\}.}
    \]

    \begin{idee}
    De façon générale, selon la constante imposée au second membre, on peut obtenir :
    \begin{itemize}
        \item aucun point si le rayon au carré obtenu est négatif ;
        \item un seul point si ce rayon au carré vaut \(0\) ;
        \item un cercle s'il est strictement positif.
    \end{itemize}
    Ici, on obtient \(3>0\), donc c'est bien un cercle.
    \end{idee}
\end{enumerate}

\vspace{1.2em}

\begin{center}
    \rule{0.65\textwidth}{0.5pt}\\[0.4em]
    \textit{Fin du corrigé}\\[0.2em]
    \rule{0.65\textwidth}{0.5pt}
\end{center}

\end{document}